% Source: Supplement: undirected low-link; https://cp-algorithms.com/graph/bridge-searching.html
% Review and validation: docs/validation.md.
\subsection{Tarjan 割点、割边与边双连通分量}
\label{algo:cuts}
\AlgorithmInfo
    {寻找无向图的割点、桥，并按去掉全部桥后的连通块划分边双分量。}
    {时间、空间 $O(n+m)$；递归 DFS，最坏深度 $O(n)$。}
    {无向边只添加一次，返回从 0 开始的边编号；允许重边、自环和孤立点。}

\subsubsection*{原理与判定条件}
割点是删除后使图的连通块数增加的顶点，桥是删除后使连通块数增加的边。对 DFS 树边 $u\to v$，$low_v$ 表示子树能通过回边到达的最早祖先：
\begin{itemize}
    \item 若 $low_v>dfn_u$，子树不能绕回 $u$ 或更早祖先，边 $(u,v)$ 是桥。
    \item 若 $u$ 不是 DFS 根，且有子节点满足 $low_v\geq dfn_u$，删除 $u$ 会隔开该子树，$u$ 是割点。
    \item DFS 根单独判断：至少有两个搜索树儿子时才是割点。不是按邻接表长度计数。
\end{itemize}
桥使用严格大于，割点使用大于等于，因为子树能回到 $u$ 本身仍不能在删除 $u$ 后脱离它连通。

\subsubsection*{重边与分量}
遍历无向边时，只跳过当前搜索树父边的\textbf{边编号}，不能跳过所有通向父顶点的边。两条平行边中，另一条可以作为回边，保证这两条边都不会被误判为桥；自环也不产生桥。

标记所有桥后，忽略桥重新遍历，得到模板采用的边双分量：任意一条非桥边的两个端点在同一分量。把每个分量缩成点，桥形成森林。割点可同时属于多个点双分量，本接口的 \InlineCode{bel} 不是点双归属，不能混用。

\subsubsection*{使用教程}
\InlineCode{add(u,v)} 内部添加双向邻接项，返回边编号。\InlineCode{build()} 后，\InlineCode{cut[u]} 判断割点、\InlineCode{bridge[id]} 判断桥，\InlineCode{bel[u]} 为去桥后的分量，\InlineCode{cnt} 为分量数，均覆盖所有连通块。
\begin{lstlisting}[style=mtycpp,numbers=none]
Cuts g;
g.init(4);
int a=g.add(1,2);
g.add(1,2),g.add(2,3),g.add(3,4);
g.build();
bool is_cut=g.cut[2],is_bridge=g.bridge[a];
int component=g.bel[1];
\end{lstlisting}

\needspace{7\baselineskip}
\subsubsection*{模板代码}
\CodeFile{code/graphs/cuts.cpp}
